% Supporting insertion, not a separate manuscript.
% Place immediately after the heading ``Triangles and their four exceptional
% markings'' in app:routing, before the historical minimal-core census.
\begin{lemma}[The reduced small affine pairs are absent from triangles]
\label{lem:current-affine-triangle-exclusions}
No actual triangular two-face has an ideal corner whose varying affine
factor and marked pair are one of the following:
\begin{enumerate}
\item a cycle of order ten, marked at antipodal nodes;
\item a cycle of order eleven, marked at distance five;
\item $\widetilde D_9$ of order ten, marked at leaves of opposite forks.
\end{enumerate}
This assertion allows arbitrary additional supporting roots and other
factors of the cusp.
\end{lemma}
\begin{proof}
Let $a,b$ be the two marked roots at the ideal corner, let $F$ be its
varying affine factor, and let $c$ be the third side root of the
triangle. Normalize every root to square two. All nodes of
$F\setminus\{a,b\}$ support the triangle. The other two actual
corner diagrams are positive semidefinite, so both principal sets
$(F\setminus\{a\})+c$ and $(F\setminus\{b\})+c$ are positive
semidefinite.

For either positive-definite deleted-factor matrix $B$, put
$p_j=-(c,r_j)\ge0$. Schur complementation gives
\[
 p^TB^{-1}p\le2.
\]
The inverse is entrywise nonnegative, and every nonzero native contact
has magnitude at least one. Consequently a position with
$(B^{-1})_{jj}>2$ cannot be contacted. This is a necessary condition
on a principal submatrix; extra supporting roots cannot repair its
failure.

Number a cycle of order $m$ by $0,\ldots,m-1$, and take $a=0,b=5$.
Deleting a node gives $A_{m-1}$, whose inverse diagonal at path
position $k$ is $k(m-k)/m$. For $m=10,11$, all positions at least
three steps from both path ends have diagonal at least
$21/10,24/11$, respectively, both greater than two. The two
deleted-factor tests therefore give these possible contact sets:
\[
\begin{array}{c|c|c}
 m&\text{after deleting }0&\text{after deleting }5\\\hline
 10&\{1,2,8,9\}&\{3,4,6,7\}\\
 11&\{1,2,9,10\}&\{3,4,6,7\}.
\end{array}
\]
In each row the sets are disjoint. Contact to either removed mark
is forbidden by the other deletion: the relevant diagonal is
$5/2$ at order ten and $30/11$ at order eleven. Thus $c$ is
orthogonal to every root of $F$.

For the last case number $F$ as the spine $0,\ldots,5$, with
leaves $6,7$ at $0$ and $8,9$ at $5$, and take $a=6,b=8$.
Each deletion is finite $D_9$. After deleting $6$, its only
inverse diagonals at most two occur at the long endpoint $7$
and its neighbor $0$, with values one and two. After deleting
$8$, they occur at $9,5$ with the same values. The short-leaf
diagonals are $9/4$, and all other diagonals exceed two.
Hence the possible contact sets are $\{0,7\}$ and $\{5,9\}$,
which are disjoint. The other deletion also forbids contact to
each removed mark, again by its short-leaf diagonal $9/4$.
Thus $c$ is orthogonal to all of $F$ in this case as well.

Let $\nu=\sum_{j\in F}\delta_jr_j$, with every $\delta_j>0$,
be the positive affine-kernel combination. It is a nonzero future
null vector representing the actual ideal corner: nonzeroness
follows by pairing with a strict interior point. The third side
of a nondegenerate actual triangle does not contain that corner,
so $-(c,\nu)>0$. The orthogonality just proved instead gives
$(c,\nu)=0$, a contradiction. Rotation and reflection cover all
cycle pairs, and fork interchange and spine reversal cover all
four opposite-fork pairs. The argument has used no restriction
on spectator rank or on the other cusp factors.
\end{proof}

We now compare the inherited triangle witnesses with the actual weights
of Appendix~\ref{app:weights}. The two cycle orbits above are the only
current reductions of ideal paths of length at most five. The reduced
order-ten $D$ orbit has distance seven and is also absent from
triangles by the lemma. Every other printed replacement in affine
$B,C,D$ of orders ten through seventeen has distance at least eight.
For larger orders the selected length-six or length-seven feature
paths retain their units: their branch or nonsimple-bond witness
persists within five steps of a mark. Such a witness cannot lie in
an affine $A$ cycle. Every affine factor of order at most nine still
pays a unit on every pair.

The four alternate witnesses below have an additional protection.
Their fixed simple leaf with a degree-two neighbor excludes every
classical affine factor, while exceptional affine factors have order
at most nine and retain full units. If the marked component is
finite, the protected leaf is at distances $2,1,0,0$ from a mark
in the four displayed cases. It remains a whole-component endpoint,
so the deep finite deletions, which require both marks to be at
least six steps from every finite feature, do not remove the unit.
The other selected finite length-six or length-seven witnesses
already have a persistent branch or nonsimple feature within five
steps. The sparse finite replacements in ranks twelve through
seventeen are all longer than these selected paths. Thus the
two-unit triangle argument below applies to the current rule,
rather than assuming that every historical short ideal pair
still has weight one.
